\documentclass[12pt]{article}
% Préambule commun à tous les documents extraits de « La Revue CAPES »
\usepackage[T1]{fontenc}
\usepackage[utf8]{inputenc}
\usepackage{lmodern}
\usepackage{amsmath,amssymb,amsthm,amsfonts,mathrsfs}
\usepackage{setspace}
\usepackage{listings}
\usepackage{pgfplots}
\pgfplotsset{compat=1.18}
\usepackage{longtable}
\usepackage{adjustbox}
\usepackage{lettrine}
\usepackage{tkz-tab}
\usepackage{changepage}
\usepackage{pdflscape}
\usepackage{graphicx}
\usepackage{tikz}
\usepackage{xcolor}
\usepackage[a4paper,top=2cm,bottom=2.5cm,left=2.5cm,right=2cm]{geometry}
\usepackage{fancyhdr}
\usepackage{draftwatermark}
\SetWatermarkText{La RevueCapes}
\SetWatermarkLightness{0.9}
\SetWatermarkScale{2}
\usepackage{hyperref}
\hypersetup{colorlinks=true,linkcolor=blue!50!black,urlcolor=blue!50!black}

\definecolor{revue}{RGB}{20,60,120}

\pagestyle{fancy}
\fancyhf{}
\renewcommand{\headrulewidth}{0.4pt}
\setlength{\headheight}{15pt}

% \entete{Matière}{Titre}{Type de document}
\newcommand{\entete}[3]{%
  \fancyhead[L]{\small\textsc{La Revue CAPES} -- #1}%
  \fancyhead[R]{\small #2}%
  \fancyfoot[C]{\thepage}%
  \thispagestyle{plain}%
  \begin{center}
    {\color{revue}\rule{\textwidth}{1.2pt}}\\[4pt]
    {\small\textsc{Concours directs CAP-CEG / CAPES -- Burkina Faso}}\\[6pt]
    {\LARGE\bfseries\color{revue} #2}\\[6pt]
    {\large\itshape #1 -- #3}\\[2pt]
    {\color{revue}\rule{\textwidth}{1.2pt}}
  \end{center}
  \vspace{0.5cm}%
}

% Encadré pour les remarques ajoutées dans les corrigés
\newcommand{\remarque}[1]{\par\medskip\noindent\fcolorbox{revue}{revue!6}{\parbox{\dimexpr\linewidth-2\fboxsep-2\fboxrule}{\textbf{\color{revue}Remarque.} #1}}\par\medskip}


\begin{document}
\entete{Culture générale}{Corrigé du sujet type de dissertation n°2}{Correction}
\begin{spacing}{1.5}

\noindent\textbf{Sujet.} Dans quelle mesure la montée du terrorisme au Burkina Faso (notamment l'action des groupes jihadistes comme le JNIM ou l'EIGS) est-elle liée aux failles de la gouvernance étatique et aux inégalités socio-économiques ?

\rubrique{1. Analyse du sujet}

\begin{itemize}
\item \textbf{Type de sujet :} question « Dans quelle mesure\dots ? ». Elle appelle un plan dialectique : montrer en quoi le lien existe, puis en marquer les limites (d'autres facteurs interviennent), avant de conclure par une appréciation nuancée et des pistes de solution.
\item \textbf{Mots-clés :}
  \begin{itemize}
  \item « montée du terrorisme » : multiplication, depuis le milieu des années 2010, des attaques contre les civils, les forces de défense et de sécurité, les écoles et les symboles de l'État ;
  \item « JNIM » (Groupe de soutien à l'islam et aux musulmans, lié à Al-Qaïda) et « EIGS » (État islamique au Grand Sahara) : principaux groupes armés opérant au Sahel ;
  \item « failles de la gouvernance » : faible présence de l'État, corruption, injustice, crise de confiance entre gouvernants et gouvernés ;
  \item « inégalités socio-économiques » : pauvreté, chômage des jeunes, écarts entre régions et entre groupes sociaux.
  \end{itemize}
\item \textbf{Problématique :} les causes internes (gouvernance défaillante et inégalités) suffisent-elles à expliquer l'expansion du terrorisme au Burkina Faso, ou faut-il aussi tenir compte de facteurs extérieurs et idéologiques ?
\end{itemize}

\rubrique{2. Plan détaillé}

\begin{enumerate}
\item[I.] \textbf{Les failles de la gouvernance : un terrain favorable au terrorisme.}
  1. L'absence ou la faiblesse de l'État dans les zones périphériques. 2. La corruption et l'injustice, sources de défiance. 3. L'instabilité politique et l'affaiblissement de l'appareil sécuritaire.
\item[II.] \textbf{Les inégalités socio-économiques : un terreau de recrutement.}
  1. La pauvreté et le chômage des jeunes. 2. Les disparités régionales et le sentiment d'abandon. 3. Les conflits fonciers et communautaires instrumentalisés.
\item[III.] \textbf{Des causes nécessaires mais pas suffisantes : la dimension régionale et idéologique.}
  1. La contagion régionale (Libye, Mali). 2. L'idéologie et les économies criminelles. 3. Conséquence : une réponse qui doit être globale (sécuritaire, politique, sociale).
\end{enumerate}

\rubrique{3. Proposition de rédaction}

\partie{Introduction}

Depuis le milieu des années 2010, le Burkina Faso, longtemps cité comme un modèle de paix et de cohésion sociale en Afrique de l'Ouest, est confronté à une grave crise sécuritaire. Les attaques de Ouagadougou en janvier 2016, puis la multiplication des violences dans le Nord, l'Est et le Centre-Nord, ont provoqué la mort de milliers de personnes, le déplacement de plus de deux millions d'habitants et la fermeture de milliers d'écoles. Ces attaques sont principalement le fait de groupes jihadistes comme le JNIM, affilié à Al-Qaïda, et l'EIGS, affilié à l'État islamique. Une question se pose alors : dans quelle mesure cette montée du terrorisme s'explique-t-elle par les failles de la gouvernance et par les inégalités socio-économiques ? Ces facteurs internes ont sans doute préparé le terrain, mais suffisent-ils à rendre compte d'un phénomène qui dépasse les frontières nationales ? Nous verrons d'abord que les défaillances de la gouvernance ont favorisé l'implantation des groupes armés, puis que les inégalités ont nourri leur recrutement, avant de montrer que d'autres facteurs, régionaux et idéologiques, doivent aussi être pris en compte.

\partie{I. Les failles de la gouvernance : un terrain favorable au terrorisme}

\souspartie{1. Un État peu présent dans les zones périphériques}
Dans de nombreuses localités du Sahel, de l'Est ou de la Boucle du Mouhoun, l'État était déjà faiblement représenté avant la crise : peu d'écoles, de centres de santé, de routes, de forces de sécurité et d'administration. Ce vide a été occupé par les groupes armés, qui se présentent comme une autorité de substitution : ils prélèvent des taxes, règlent des conflits, imposent leurs règles. Le groupe Ansarul Islam, apparu dans la province du Soum autour de 2016, a ainsi prospéré dans des zones où l'État était perçu comme lointain.

\souspartie{2. La corruption et l'injustice, sources de défiance}
Lorsque la justice est lente ou partiale, que des fonds publics sont détournés et que certains agents abusent de leur pouvoir, les citoyens perdent confiance dans les institutions. Les groupes jihadistes exploitent ce ressentiment en dénonçant la corruption des élites et en promettant une justice expéditive mais perçue comme « équitable ». Les accusations d'exactions contre des civils commises au nom de la lutte antiterroriste ont, elles aussi, poussé certaines populations à se méfier de l'État.

\souspartie{3. L'instabilité politique et l'affaiblissement de l'appareil sécuritaire}
La chute du régime de Blaise Compaoré en octobre 2014, la tentative de coup d'État de septembre 2015 et la dissolution du Régiment de sécurité présidentielle ont désorganisé l'appareil sécuritaire et le renseignement au moment même où la menace s'étendait depuis le Mali. Les coups d'État de 2022 témoignent aussi de la difficulté des institutions à répondre efficacement à la crise.

\transition{Les failles de la gouvernance ont ainsi ouvert des brèches. Mais le terrorisme a aussi besoin de combattants : c'est là qu'interviennent les inégalités socio-économiques.}

\partie{II. Les inégalités socio-économiques : un terreau de recrutement}

\souspartie{1. La pauvreté et le chômage des jeunes}
Le Burkina Faso a une population très jeune, et beaucoup de jeunes, surtout en milieu rural, n'ont ni formation ni emploi. Les groupes armés offrent de l'argent, une moto, une arme, un statut social et un sentiment d'appartenance. Pour certains jeunes désœuvrés, l'engagement relève moins de la conviction religieuse que de la recherche de revenus, de prestige ou de protection.

\souspartie{2. Les disparités régionales et le sentiment d'abandon}
Les régions périphériques accusent souvent un retard en matière d'infrastructures, de scolarisation et de services de santé par rapport aux grands centres urbains. Ce déséquilibre nourrit un sentiment de marginalisation que les discours extrémistes transforment en colère contre l'État.

\souspartie{3. Les conflits fonciers et communautaires instrumentalisés}
La raréfaction des terres et des pâturages, aggravée par la croissance démographique et les changements climatiques, accentue les tensions entre agriculteurs et éleveurs. Les groupes armés attisent ces conflits et les présentent comme des affrontements communautaires. La stigmatisation de certaines communautés, accusées en bloc de complicité, favorise à son tour le recrutement.

\transition{Gouvernance défaillante et inégalités constituent donc des causes profondes. Toutefois, elles n'expliquent pas tout : d'autres pays connaissent la pauvreté sans subir une telle vague terroriste.}

\partie{III. Des causes nécessaires mais pas suffisantes}

\souspartie{1. La contagion régionale}
Le terrorisme au Burkina Faso s'inscrit dans une crise sahélienne plus large. La chute du régime de Kadhafi en Libye (2011) a libéré des armes et des combattants ; la crise du Nord du Mali (2012) a permis l'installation durable de groupes jihadistes, qui se sont ensuite étendus vers le Burkina Faso et le Niger. Les frontières poreuses facilitent les mouvements de combattants et d'armes.

\souspartie{2. L'idéologie et les économies criminelles}
L'idéologie jihadiste, diffusée par des prédicateurs radicaux et les réseaux sociaux, fournit une justification religieuse à la violence. Les groupes armés se financent aussi par des activités criminelles : trafics, enlèvements, vol de bétail, contrôle de sites d'orpaillage. Le terrorisme est donc aussi une entreprise politique et économique.

\souspartie{3. Une réponse nécessairement globale}
Si la gouvernance et les inégalités ne sont pas les seules causes, elles restent des leviers essentiels d'action. La réponse militaire est indispensable pour protéger les populations, mais elle doit s'accompagner du retour de l'État et des services sociaux de base, de la lutte contre la corruption, de la création d'emplois pour les jeunes, de la justice pour tous et du dialogue intercommunautaire, ainsi que d'une coopération régionale renforcée.

\partie{Conclusion}

La montée du terrorisme au Burkina Faso est largement liée aux failles de la gouvernance, qui ont laissé des territoires sans État, et aux inégalités socio-économiques, qui ont fourni aux groupes jihadistes un vivier de recrues. Ces causes internes ne suffisent cependant pas : la déstabilisation régionale, l'idéologie extrémiste et les économies criminelles ont aussi joué un rôle déterminant. La lutte contre le terrorisme ne peut donc être seulement militaire : elle suppose de reconstruire la confiance entre l'État et les citoyens et d'offrir un avenir à la jeunesse. La question reste posée de savoir comment concilier l'urgence sécuritaire et ces réformes de long terme.

\end{spacing}
\end{document}
